%%=============================================================================
%% Conclusie
%%=============================================================================

\chapter{Conclusie}%
\label{ch:conclusie}

% TODO: Trek een duidelijke conclusie, in de vorm van een antwoord op de
% onderzoeksvra(a)g(en). Wat was jouw bijdrage aan het onderzoeksdomein en
% hoe biedt dit meerwaarde aan het vakgebied/doelgroep? 
% Reflecteer kritisch over het resultaat. In Engelse teksten wordt deze sectie
% ``Discussion'' genoemd. Had je deze uitkomst verwacht? Zijn er zaken die nog
% niet duidelijk zijn?
% Heeft het onderzoek geleid tot nieuwe vragen die uitnodigen tot verder 
%onderzoek?

Dit onderzoek had als doel na te gaan hoe apps voor stress- en burn-outpreventie het mentale welzijn van Vlaamse twintigers kunnen ondersteunen en hoe inzichten uit bestaande toepassingen kunnen worden vertaald naar concrete ontwerpprincipes en functionaliteiten voor een \acrshort{poc}. Op basis van de literatuurstudie, het expertinterview, de requirementsanalyse, de analyse van bestaande toepassingen en de ontwikkeling en evaluatie van de \acrshort{poc} kunnen de onderzoeksvragen als volgt worden beantwoord.

\vspace{1em}

\section{Het probleem}
\label{sec:het-probleem-conclusie}

\vspace{1em}

Vlaamse twintigers vormen geen homogene doelgroep en hebben bijgevolg ook niet allemaal dezelfde noden op vlak van mentale gezondheid. Binnen generatie Z kan voornamelijk een onderscheid worden gemaakt tussen studenten en jongvolwassenen die actief zijn op de arbeidsmarkt. Studenten ervaren onder meer academische verwachtingen en prestatiedruk, terwijl werkende twintigers bijkomende druk kunnen ervaren door werkbelasting, loopbaanverwachtingen en financiële onzekerheid. Tegelijkertijd zijn er verschillende gedeelde kenmerken, zoals perfectionisme, ambitie, sociale vergelijking en een sterke behoefte aan erkenning en persoonlijke ontwikkeling. Vooral de psychologische basisbehoeften aan autonomie, competentie en verbondenheid zijn hierbij relevant. Een ondersteunende toepassing moet daarom voldoende flexibel zijn om rekening te houden met verschillende levenssituaties, terwijl ze tegelijkertijd deze gedeelde behoeften ondersteunt.

\newpage

De manier waarop deze persoonlijke en externe factoren worden ervaren, heeft ook invloed op de gebruikerservaring van digitale toepassingen. Werk-, studie- en sociale druk kunnen ervoor zorgen dat gebruikers reeds met een hoge mentale belasting een toepassing openen. Een hoge mate van digitale connectiviteit, voortdurende bereikbaarheid en sociale vergelijking kunnen deze belasting verder versterken door information overload, technostress en een verhoogde alertheid voor digitale prikkels. Een toepassing die mentale ondersteuning wil bieden, moet daarom vermijden dat ze zelf bijkomende inspanning vraagt. Duidelijke visuele hiërarchie, beperkte informatie per scherm, korte teksten, eenvoudige navigatie en het opdelen van complexe informatie in kleinere stappen zijn hierdoor belangrijke voorwaarden voor een ondersteunende gebruikerservaring.

\vspace{1em}

Het bestaande aanbod voor mentale gezondheid biedt hiervoor nog geen volledige oplossing. Mentale gezondheidsapps richten zich voornamelijk op mindfulness, meditatie en ademhalingsoefeningen, terwijl andere interventievormen, zoals \acrshort{rebt}, minder vertegenwoordigd zijn. Ook toepassingen die specifiek gericht zijn op burn-outpreventie blijken schaars en er is weinig sprake van één geïntegreerde toepassing die verschillende preventieve interventies combineert. Daarnaast kunnen complexe navigatie, beperkte personalisatie, overvolle interfaces, onduidelijke communicatie en een grote hoeveelheid notificaties de cognitieve belasting verhogen. Het falen van bestaande toepassingen ligt daardoor niet noodzakelijk in een gebrek aan inhoudelijke interventies, maar vooral in de combinatie van inhoud, gebruikerservaring en de context waarin deze wordt aangeboden.

\vspace{1em}

Binnen de bestaande apps voor mentaal welzijn worden voornamelijk mindfulness, meditatie, ademhalingsoefeningen, dagboeken, stemmingsregistratie en psycho-educatie toegepast. Deze functionaliteiten kunnen waardevol zijn, maar hun meerwaarde wordt mede bepaald door de manier waarop ze worden aangeboden. De analyse wijst daarom op het belang van aanvullende principes zoals personalisatie, autonomie, overzichtelijkheid, beperkte digitale prikkels en ondersteuning van persoonlijke vooruitgang. Een goede mentale gezondheidsapp hoeft bijgevolg niet zoveel mogelijk functionaliteiten te bevatten, maar moet vooral de juiste functionaliteiten op een samenhangende en laagdrempelige manier aanbieden.

\vspace{1em}


\section{Functionele en niet-functionele vereisten}
\label{sec:requirements-conclusie}

\vspace{1em}

Voor een toepassing die stress en burn-out bij twintigers helpt voorkomen, zijn zowel inhoudelijke functionaliteiten als kwaliteitsvoorwaarden essentieel. Tot de belangrijkste functionaliteiten behoren onboarding en personalisatie, stressregistratie en zelfreflectie, een gestructureerd dagboek met historiek, psycho-educatie, stressreducerende oefeningen, feedback en zichtbare persoonlijke vooruitgang. 

Ook instelbare notificaties en mogelijkheden tot ondersteuning en doorverwijzing naar professionele hulp zijn relevant. Deze functionaliteiten moeten worden aangeboden binnen kwaliteitsvoorwaarden zoals bruikbaarheid, toegankelijkheid, performantie, betrouwbaarheid, beveiliging en privacy. Zeker bij gevoelige mentale gezondheidsinformatie moeten gebruikers controle behouden over hun gegevens en over de manier waarop zij met de toepassing interageren.

\vspace{1em}


\section{UI/UX en mentale rust}
\label{sec:UI-UX-mentale-rust}

\vspace{1em}

Een rustige gebruikerservaring ontstaat vooral door informatie en prikkels bewust te beperken. Overzichtelijkheid, een duidelijke visuele hiërarchie, consistente interacties, beperkte informatie per scherm en eenvoudige navigatie dragen bij aan mentale rust. Ook chunking, progressive disclosure, korte scanbare teksten en één primaire taak per scherm beperken de hoeveelheid informatie die gebruikers tegelijk moeten verwerken. Een rustig kleurenpalet, beperkt gebruik van fonts en kleuren, weinig animaties en terughoudend gebruik van afbeeldingen en video helpen daarnaast visuele overprikkeling te voorkomen. Complexe navigatie, overvolle interfaces, inconsistente interacties en grote hoeveelheden informatie werken daarentegen in de tegenovergestelde richting en verhogen de cognitieve belasting.

\vspace{1em}

Deze inzichten kunnen worden vertaald naar concrete ontwerpprincipes voor een interface die mentale rust bevordert. Daarbij staan onder meer privacy by design, inclusieve toegankelijkheid, beperkte digitale interrupties, behoud van overzicht, minimalisering van cognitieve belasting, het vermijden van visuele overprikkeling en techno-simplicity centraal. Daarnaast moeten ontwerpkeuzes autonomie, competentie, verbondenheid, persoonlijke relevantie en betekenisvolle vooruitgang ondersteunen. De interface moet de gebruiker dus niet alleen informeren of activeren, maar vooral het gevoel geven dat de toepassing ondersteunend, beheersbaar en relevant is.

\vspace{1em}

Ook notificaties en gamification vragen om een terughoudende aanpak. Meldingen kunnen helpen bij gedragsverandering en engagement, maar kunnen tegelijkertijd digitale overbelasting en online vigilance versterken. Ze moeten daarom optioneel, instelbaar, beperkt, contextafhankelijk en gepersonaliseerd zijn en mogen niet dwingend of urgent worden geformuleerd. Gamification kan eveneens ondersteuning bieden, zolang de nadruk ligt op persoonlijke groei en betekenisvolle vooruitgang. Punten, rankings en competitieve streaks zijn minder geschikt omdat ze de focus kunnen verschuiven van welzijn naar prestatie. Spelelementen moeten bijgevolg het gezondheidsdoel ondersteunen en geen nieuwe prestatiedruk creëren.

\vspace{1em}


\section{Beantwoording van de hoofdonderzoeksvraag}
\label{sec:beantwoording-hoofdonderzoeksvraag}

\vspace{1em}

Apps voor stress- en burn-outpreventie kunnen het mentale welzijn van Vlaamse twintigers ondersteunen wanneer inhoudelijke interventies worden gecombineerd met een gebruikerservaring die zelf geen bijkomende bron van stress vormt. De resultaten tonen aan dat hiervoor niet alleen relevante functionaliteiten nodig zijn, maar ook een sterke focus op personalisatie, autonomie, overzichtelijkheid, beperkte digitale prikkels en cognitieve rust. De opgestelde requirements en ontwerpprincipes bieden hiervoor een concreet kader, terwijl de \acrshort{poc} aantoont dat deze inzichten kunnen worden vertaald naar concrete ontwerp- en ontwikkelkeuzes. De meerwaarde ligt daarmee niet in het aanbieden van zoveel mogelijk functies, maar in het doelgericht combineren van functionaliteit, interface en interactie zodat de toepassing ondersteunt zonder bijkomende cognitieve belasting of prestatiedruk te veroorzaken.

\vspace{1em}

De resultaten liggen grotendeels in lijn met de verwachtingen uit het oorspronkelijke onderzoeksvoorstel. Zo werd bevestigd dat mindfulness, meditatie en ademhalingsoefeningen sterk vertegenwoordigd zijn in bestaande toepassingen, terwijl andere interventievormen minder vaak voorkomen. Ook kwamen verschillende potentiële bronnen van cognitieve belasting, zoals complexe navigatie, grote hoeveelheden informatie en frequente notificaties, terug in de analyse. Dit bevestigt dat niet alleen de inhoud van een mentale gezondheidsapp, maar ook de manier waarop deze wordt aangeboden, bepalend is voor de potentiële meerwaarde ervan.

\vspace{1em}


\section{Beperkingen}
\label{sec:beperkingen}

\vspace{1em}

Hoewel deze bachelorproef een onderbouwde ontwerpaanpak en \acrlong{poc} heeft opgeleverd, zijn er verschillende beperkingen. De belangrijkste beperking is dat de ontwikkelde \acrshort{poc} niet werd geëvalueerd met eindgebruikers. De evaluatie gebeurde voornamelijk aan de hand van een heuristische beoordeling van de requirements en ontwerpprincipes, aangevuld met een evaluatie door de copromotor. Hierdoor kon worden nagegaan in welke mate de voorgestelde ontwerpkeuzes werden toegepast, maar niet of gebruikers deze daadwerkelijk als rustgevend, begrijpelijk of ondersteunend ervaren. De resultaten tonen daarom vooral aan dat de voorgestelde ontwerpaanpak praktisch toepasbaar is en vormen geen bewijs dat de toepassing daadwerkelijk stress- of burn-outklachten vermindert.

\vspace{1em}

Een tweede beperking betreft de concretisering van de ontwerpprincipes. Het principe rond inclusieve toegankelijkheid is bijvoorbeeld relatief breed geformuleerd en verwijst naar de \acrshort{wcag}-richtlijnen, die een omvangrijk geheel aan criteria bevatten. Binnen de scope van deze bachelorproef was het niet haalbaar om deze volledig te vertalen naar concrete en controleerbare criteria. Daarnaast werden de ontwerpprincipes door de onderzoeker zelf toegepast. De aanwezige voorkennis uit de opleiding grafisch ontwerp en de verdieping in het onderzoeksdomein kunnen daardoor invloed hebben gehad op de interpretatie en toepassing ervan. Het is bijgevolg nog niet duidelijk in welke mate andere ontwerpers of ontwikkelaars dezelfde principes op een gelijkaardige manier zouden interpreteren.

\vspace{1em}

Ten slotte is de toepasbaarheid van de ontwerpprincipes buiten het domein van mentale gezondheidsapps nog onvoldoende onderzocht. Hoewel verschillende principes, zoals het beperken van cognitieve belasting, het ondersteunen van autonomie en het vermijden van prestatiedruk, mogelijk ook relevant zijn voor andere digitale toepassingen, werd dit binnen deze bachelorproef niet onderzocht.

\vspace{1em}

\section{Aanbevelingen voor verder onderzoek}
\label{sec:verder-onderzoek}

\vspace{1em}

Een eerste belangrijke stap voor verder onderzoek is een evaluatie met eindgebruikers. Een gebruikerstest en, bij voorkeur, een langdurige evaluatie kunnen nagaan of gebruikers de toepassing daadwerkelijk als minder cognitief belastend ervaren en of de ontwerpkeuzes op langere termijn samenhangen met veranderingen in ervaren stress. Hierbij kan zowel worden onderzocht of gebruikers de functionaliteiten begrijpen en waarderen als of de ontwerpprincipes daadwerkelijk bijdragen aan een ondersteunende gebruikerservaring.

\vspace{1em}

Daarnaast is verder onderzoek nodig naar de overdraagbaarheid en reproduceerbaarheid van de ontwerpprincipes. Omdat de principes binnen deze bachelorproef door de onderzoeker zelf werden toegepast, kan toekomstig onderzoek waarbij externe ontwerpers of developers de richtlijnen toepassen inzicht geven in de mate waarin deze voldoende concreet en eenduidig zijn. Dit kan bovendien aantonen welke principes verdere concretisering nodig hebben.

\vspace{1em}

Tot slot kan worden onderzocht in welke mate de ontwerpprincipes toepasbaar zijn buiten mentale gezondheidsapps. Principes zoals het beperken van cognitieve belasting, het vermijden van prestatiedruk, het respecteren van rustmomenten en het ondersteunen van autonomie zijn niet noodzakelijk gebonden aan mentale gezondheid. Het kan daarom waardevol zijn om na te gaan of dezelfde ontwerpaanpak ook kan bijdragen aan een minder belastende gebruikerservaring binnen bijvoorbeeld sociale media, productiviteitsapplicaties of andere digitale toepassingen waarin gebruikers grote hoeveelheden informatie en taken moeten verwerken.